\newpage
% phantomsection: eigener hyperref-Anker, sonst zeigt der TOC-Link auf das
% vorherige Verzeichnis statt auf diese Seite.
\phantomsection
\addcontentsline{toc}{section}{\langde{KI-Verzeichnis}\langen{List of AI tools}}
\section*{\langde{KI-Verzeichnis}\langen{List of AI tools}}

% Dieses Verzeichnis weist die im Rahmen der Arbeit eingesetzten Werkzeuge
% Künstlicher Intelligenz aus (gemaess dem FOM-Leitfaden zur Gestaltung
% wissenschaftlicher Arbeiten). Bitte an die tatsaechlich genutzten Werkzeuge
% anpassen oder das Verzeichnis entfernen, falls keine KI-Werkzeuge genutzt
% wurden. Art und Umfang der Nutzung werden ueblicherweise zusaetzlich im
% Methodik-Kapitel beschrieben; die "Erklärung zur KI-Nutzung" befindet sich
% im Anhang (siehe kapitel/anhang/erklaerung.tex).

\langde{Das folgende Verzeichnis weist die im Rahmen dieser Arbeit eingesetzten
Werkzeuge Künstlicher Intelligenz aus. Art und Umfang der Nutzung sind im
Methodikteil der Arbeit beschrieben.}
\langen{The following list documents the artificial-intelligence tools used in
the course of this work. The nature and extent of their use are described in the
methodology part of the work.}

\begin{description}
  % Beispieleintrag -- bitte anpassen bzw. entfernen:
  \item[Anbieter (Jahr)] \langde{Produktname}\langen{Product name}
  [\langde{kurze Beschreibung des Werkzeugs und des Sprachmodells}%
   \langen{short description of the tool and the underlying model}],
  \langde{Abfragen Monat/Jahr}\langen{queried month/year}.
  \url{https://example.org}
\end{description}
